\chapter{Introdução}
\label{chp:introducao}

A metodologia expositiva de ensino, iniciada com o patriarca da pedagogia, Platão (427–347 a.C.), é uma técnica tradicional, que perdura até os dias atuais \cite{andreataAulaExpositivaPaulo2019}; apesar dos inúmeros avanços tecnológicos ocorridos ao longo do tempo. Seu elemento central é a oratória, por meio da qual o educador transmite seu saber aos alunos. Bartolomeis(1984, \textit{apud}) \cite{andreataAulaExpositivaPaulo2019} descreve a prática da seguinte forma: "Acusamos de verbalismo a Escola Tradicional, onde, no entanto, é  sobretudo o professor quem fala, uma vez que os alunos só o devem  fazer quando interrogados. Os seus estados normais são, por isso, a imobilidade e o silêncio.". Dito isso, evidencia-se uma hierarquia clara: o professor ocupa o centro do processo de aprendizagem, enquanto o aluno assume o papel de ouvinte, um coadjuvante, sendo responsável pela retenção do conteúdo.

No entanto, ainda que tradicional, a prática deve ser revista, considerando o atual contexto tecnológico. Autores como \cite{andreataAulaExpositivaPaulo2019} e \cite{pereiraCriticaMetodologiaTradicional2022} demonstram preocupação com a retenção do conteúdo pelo aluno; pois esse tem uma papel passivo em seu próprio processo de aprendizado. Além disso, hoje, é preciso se preocupar também com os inúmeros estímulos externos advindos da exposição às telas \cite{lopesUsoCelularEm2017}. Essa realidade torna-se ainda mais preocupante em países com dificuldades de acesso à tecnologia, como o Brasil, onde cresce o número de estudantes no Ensino Superior em modalidades à distância (EAD) \cite{moranEnsinoSuperiorDistancia2009}, sem que haja necessariamente um aumento no acompanhamento individualizado desses alunos, para garantir que estão aprendendo.

\section{Problema}
A metodologia expositiva tem origem no Mundo Antigo \cite{andreataAulaExpositivaPaulo2019}. Nesse sentido, é evidente que perseverou por incontáveis séculos, atravessando diferentes paradigmas tecnológicos e perdurando até o contexto atual. É uma técnica na qual o emissor é protagonista no processo de ensino, enquanto o ouvinte --- que devia receber maior atenção para a retenção do conteúdo --- é relegado a coadjuvante nesse vital processo.

Essa problemática é constantemente levantada por acadêmicos, como é o caso de \citeonline{vasconcellosMetodologiaDialeticaEm1992} e \citeonline{pereiraCriticaMetodologiaTradicional2022}, que apontam o mero papel de ouvinte como prejudicial aos estudantes, que não são estimulados a engajar; tem uma postura meramente passiva. Eles ainda citam que o docente procura verificar se todos estão entendendo, mas por meio de perguntas vagas, indagando os estudantes quanto à sua compreensão. Nesse processo, os estudantes apenas concordam, e a conversa é finalizada no ato. Não há acompanhamento ativo, apenas uma estrutura hierarquizada e de efetividade discutível.

Os estudos de \citeonline{vasconcellosMetodologiaDialeticaEm1992} e \citeonline{pereiraCriticaMetodologiaTradicional2022} foram feitos em ambiente presencial de ensino, no  qual há mais facilidade de interação entre as partes, além de melhor controle sobre a atenção e engajamento da turma. Então, se mesmo nessas condições ambos os estudos se mostram céticos quanto à efetividade da aula expositiva, é de se esperar que em ambiente remoto esse cenário seja ainda pior.

Com a pandemia de COVID-19, o ensino a distância ganhou ainda mais espaço, especialmente no Ensino Superior, que já apresentava tendência de crescimento exponencial no Brasil \cite{moranEnsinoSuperiorDistancia2009}. Apesar de ser um país relativamente pobre e com acesso limitado à tecnologia, o Brasil registra uma proporção significativa de estudantes universitários matriculados em cursos a distância, tanto totalmente online quanto semipresenciais.

\section{Objetivos}
Este trabalho busca investigar o desempenho da metodologia expostivia de ensino, aplicada ao ensino remoto; tendo a análise pautada em indicadores de foco e engajamento, emitidos em relatórios pela extensão ESPEON.

\subsection{Objetivos Específicos}
Para viabilizar a investigação do cenário proposto, estabelece-se como objetivo específico a construção do \textit{software} ESPEON, uma extensão para navegadores de \textit{internet}. O \textit{software} utiliza \textit{snapshots} das guias do navegador para, com base nas propriedades das abas, construir relatórios que avaliam o comportamento discente durante aulas expositivas. A ferramenta é empregada na coleta de dados comportamentais de usuários submetidos a aulas expositivas em ambiente remoto, possibilitando a identificação de padrões relacionados à dispersão de foco e à participação.

\section{Justificativa}\label{sec:justificativa}
O avanço das tecnologias digitais e a popularização do ensino remoto transformaram profundamente as práticas educacionais, ampliando o acesso ao conhecimento, mas também impondo novos desafios aos professores. Entre esses desafios, destacam-se a dificuldade em manter a atenção e o engajamento dos estudantes, bem como a limitação de instrumentos que permitam observar e mensurar esses aspectos de forma sistemática.

Em ambientes virtuais de aprendizagem, a ausência do contato presencial e a variedade de estímulos externos podem favorecer a dispersão e a perda de foco, comprometendo o rendimento acadêmico e a efetividade das metodologias de ensino. Nesse contexto, torna-se importante fazer um acompanhamento ativo do discente e como ele interage nesse ambiente virtual, se faz sentido seguir com a aula expositiva nesse âmbito.

O \textit{software} ESPEON surge como uma proposta para atender a essa necessidade, monitorar o comportamento dos alunos em sala de aula virtual. A partir de métricas de tempo e estado das abas do navegador, é possível identificar padrões de engajamento e desatenção, fornecendo subsídios para aprimorar estratégias pedagógicas e compreender melhor a dinâmica da aprendizagem mediada por tecnologia.

Dessa forma, o presente estudo justifica-se pela relevância de se investigar o desempenho da metodologia expositiva no ensino remoto a partir de dados empíricos coletados em tempo real. Além de contribuir para a discussão sobre atenção e engajamento em ambientes virtuais, o trabalho propõe uma solução tecnológica de baixo custo e alta aplicabilidade, consolidando um instrumento útil para pesquisas futuras e práticas docentes.
\section{Contribuições}

Esta pesquisa contribui para a compreensão da interatividade em ambientes de sala de aula virtual, tendo como parâmetro o contexto da metodologia expositiva, e servindo como norteadora de tendências para esse cenário. Por meio do ESPEON, essa avaliação torna-se acessível e prática, em razão da baixa barreira de entrada para uso e implementação da ferramenta. Além disso, o monitoramento ocorre de forma transparente, em \textit{background}, minimizando desconfortos ao usuário durante o processo. 

Os relatórios de aula desempenham papel central na pesquisa, por sua flexibilidade de formato — podendo se apresentar tanto como um documento PDF simples quanto como relatórios mais elaborados, gerados a partir do consumo da \textit{API} de relatórios do sistema. Essa versatilidade facilita a condução de estudos futuros e a consolidação de dados provenientes de diferentes contextos, possibilitando análises mais abrangentes e precisas sobre o fenômeno investigado.

\section{Questão de Pesquisa (QP)}\label{sec:research.question}

A questão de pesquisa tem como finalidade orientar o desenvolvimento deste estudo, direcionando a análise sobre a efetividade da metodologia expositiva quando aplicada ao contexto do ensino remoto. Dessa forma, constitui-se como o elemento norteador da monografia, formulado a partir do objetivo principal do trabalho.

\textbf{QP:} De que maneira a metodologia tradicional de ensino expositivo se manifesta em sala de aula virtual, do ponto de vista do foco e engajamento dos estudantes?

\section{Hipótese}
Parte-se da hipótese de que alunos submetidos à metodologia expositiva em sala de aula virtual tendem a apresentar dispersão de foco e baixo engajamento, em razão das múltiplas distrações presentes em seu ambiente doméstico e da limitação do docente em monitorar e intervir sobre comportamentos que ocorrem de forma velada por trás da tela.

A formulação de um problema solucionável é fundamental para o início de uma pesquisa científica, sendo a etapa seguinte a proposição de uma solução, denominada hipótese <citar-gil>.

Para este trabalho, propõe-se uma abordagem de associação entre variáveis, na qual as métricas coletadas serão associadas a variáveis representativas de atenção e engajamento. Serão observados, respectivamente, o tempo fora da conferência (aula) e o estado dos periféricos. <citar-gil> afirma que esse tipo de hipótese busca indicar a força ou o sentido da relação variável-valor, mas não determina causa, dependência ou influência.
% \section{Referencial teórico}
\begin{enumerate}
    \item\label{rt_migrations} Migrações - as migrações são classes que representam interações pontuais com a base de dados, geralmente utilizadas para versionamento da estrutura da base de dados ou para alterações pontuais nos dados quando necessário.
    \item\label{rt_orm} Mapeador objeto-relacional (ORM) - os mapeadores objeto-relacional são abstrações que tem como premissa servir de ponte para duas linguagens, no caso do Vortex seriam o PHP que é uma linguagem orientada a objetos e o SQL que é a linguagem da base de dados. O mapeador age traduzindo o contexto de um lado da ponte para outro e vice-versa.
    \item\label{rt_reflections} \textit{Reflections} -   definida por \cite{gabriel1991clos}, reflexão é a capacidade de um programa manipular como dados algo que representa o estado do programa durante sua própria execução. Há dois elementos nessa manipulação: introspecção e intercessão. A introspecção refere-se à capacidade de um programa observar e, consequentemente, raciocinar sobre seu próprio estado. A intercessão é a habilidade de um programa modificar seu próprio estado de execução ou alterar sua própria interpretação ou significado. Ambos os aspectos requerem um mecanismo para codificar o estado de execução como dados; esse processo é chamado de reificação.
    \item\label{rt_framework} \textit{Framework} - conjunto de abstrações e implementações que visam facilitar ou incrementar a produção de determinada aplicação.
    \item\label{rt_guessers} \textit{Guessers} - são objetos ou métodos relacionados à ''adivinhação'' ou estimativas com base em lógicas predefinidas.
\end{enumerate}
\section{Estrutura do trabalho}
Este estudo está estruturado em mais cinco capítulos distintos. No capítulo \ref{sec:estado_atual}, apresenta-se em detalhes a metodologia abrangente que norteou o desenvolvimento do Vortex, abordando desde sua fundação tecnológica até suas políticas, padrões de projeto e boas práticas de desenvolvimento seguidas. Nessa seção, são delineados os principais conceitos, funcionalidades e classes que compõem integralmente o referido \textit{framework}.

No capítulo \ref{sec:modelagem} estão descritas as classes e suas responsabilidades dentro do \textit{framework}, assim como suas relações com as demais, além das representações dos diagramas das mesmas.

O subsequente capítulo \ref{Utilização} fornece uma descrição pormenorizada do procedimento adequado para a utilização de cada funcionalidade do Vortex. Tal exposição abrange desde comandos do Cosmo até as classes que integram a arquitetura global do \textit{framework}.

O capítulo \ref{sec:estudo_de_caso} é dedicada ao estudo de caso, onde são delineados os testes conduzidos para verificar a superioridade de desempenho do Vortex em comparação com outros \textit{frameworks} similares. Essa comprovação é estabelecida mediante a realização de experimentos controlados em cenários equivalentes.

Finalmente, a conclusão, abordada no capítulo \ref{chp:conclusao}, destaca os resultados obtidos ao longo do estudo, proporcionando \textit{insights} conclusivos. Além disso, são delineadas as perspectivas para trabalhos futuros, consolidando assim a contribuição e relevância do presente trabalho.